% ============================================================
% PHẦN 45 phút — Slide 38-41 (Quan sát RAM, Tổng kết pipeline,
% Tổng kết FastGS & giới hạn, Cảm ơn / Q&A) — rút gọn từ part7.tex
% ============================================================

% --- Slide 38 (gốc: Slide 62 — Quan sát 2) ---
\begin{frame}{Quan sát 2: Nút thắt là RAM hệ thống, không phải VRAM}
\begin{itemize}
    \item Đo trên Google Colab T4 trong quá trình huấn luyện:
\end{itemize}
\footnotesize
\begin{table}[]
\centering
\begin{tabular}{lccc}
\toprule
Tài nguyên & Đỉnh sử dụng & Khả dụng & \% sử dụng \\
\midrule
VRAM (GPU) & 1.16 GB & 14.56 GB & \textbf{8\%} \\
RAM hệ thống & 9.59 GB & 12.7 GB & \textbf{91\%} (giới hạn mềm 10.5 GB) \\
\bottomrule
\end{tabular}
\end{table}
\normalsize
\begin{itemize}
    \item RAM hệ thống áp sát giới hạn mềm, trong khi VRAM còn dư rất nhiều.
    \item Khuyến nghị cũ "giảm độ phân giải để tiết kiệm VRAM" \textbf{không cần thiết} ở ngân sách này.
\end{itemize}
\end{frame}

% --- Slide 39 (gốc: Slide 64 — Tổng kết pipeline) ---
\begin{frame}{Tổng kết pipeline render 3D Gaussian Splatting}
\centering
\footnotesize
\begin{tikzpicture}[
    node distance=8mm and 6mm,
    box/.style={draw, rounded corners, align=center, minimum height=0.9cm, minimum width=2.6cm, fill=blue!8}
]
\node[box] (gauss) {3D Gaussian\\$(\mu,\Sigma,\alpha,SH)$};
\node[box, right=of gauss] (proj) {Chiếu phối cảnh\\(EWA)};
\node[box, right=of proj] (raster) {Rasterizer khả vi\\theo tile};

\node[box, below=of raster] (blend) {Alpha\\blending};
\node[box, left=of blend] (loss) {Loss\\($L_1$+D-SSIM)};
\node[box, left=of loss] (back) {Backward\\(gradient)};

\node[box, below=of back, fill=orange!12] (adc) {Adaptive\\Density Control};

\draw[-Latex] (gauss) -- (proj);
\draw[-Latex] (proj) -- (raster);
\draw[-Latex] (raster) -- (blend);
\draw[-Latex] (blend) -- (loss);
\draw[-Latex] (loss) -- (back);
\draw[-Latex] (back) -- (adc);
\draw[-Latex] (adc) -- (gauss);
\end{tikzpicture}
\vspace{2mm}

\normalsize
Một chu trình khép kín: render $\rightarrow$ tính loss $\rightarrow$ lan truyền ngược $\rightarrow$ điều chỉnh mật độ Gaussian $\rightarrow$ lặp lại.
\end{frame}

% --- Slide 40 (gốc: Slide 65 — Tổng kết cải tiến FastGS \& giới hạn) ---
\begin{frame}{Tổng kết cải tiến FastGS \& giới hạn còn lại}
\textbf{Ba đòn bẩy chính:}
\begin{itemize}
    \item \textbf{Multi-view score} + \textbf{Compact box} $\Rightarrow$ giảm Gaussian dư thừa mà không giảm chất lượng đáng kể.
    \item \textbf{Sparse Adam} $\Rightarrow$ giảm chi phí tính toán mỗi iteration.
\end{itemize}
\vspace{2mm}
\textbf{Giới hạn / hướng phát triển:}
\begin{enumerate}
    \item Dữ liệu Gaussian đầu ra chưa được nén (248 B/Gaussian, thô).
    \item Ngân sách 7k iteration bỏ lỡ bước \texttt{final\_prune}.
    \item Chưa có ablation study cô lập riêng từng đòn bẩy.
    \item 4 nhánh mở rộng của FastGS gốc — dynamic scenes, sparse view, surface reconstruction, SLAM — chưa tích hợp vào fork này.
\end{enumerate}
\end{frame}

% --- Slide 41 (gốc: Slide 67 — Cảm ơn / Q&A, frame cuối cùng) ---
\begin{frame}[c]{}
\centering
{\Huge \textbf{Cảm ơn đã theo dõi!}}

\vspace{6mm}
{\LARGE Hỏi đáp (Q\&A)}

\vspace{10mm}
{\small Digital Twin GS --- FastGS-lite \\ \texttt{github.com/KietAnhCS/fastgs-lite}}
\end{frame}
